\uuid{fRMr}
\titre{Changement de variables}
\auteur{hanani ; mzari}
\organisation{exo7}
\licence{CC BY-NC-SA 4.0}
\datecreate{2020-01-20}
\niveau{L1}
\module{Analyse}
\chapitre{Calcul d'intégrales}
\sousChapitre{Primitives diverses}
\difficulte{2}
\theme{Primitives, changement de variables}
\exo7qcmid{lille-2.primitives.026}

\contenu{
\begin{question}
\qclassification{127.11}
\qtags{motcle=changement de variables}
Parmi les assertions suivantes, cocher celles qui sont vraies :
\begin{answers}
\bad{Le changement de variable $u=\cos x$ donne $\displaystyle \int \cos ^3x\mathrm{d}x=\int u^3\mathrm{d}u$.}
\good{Le changement de variable $u=\sin x$ donne $\displaystyle \int \cos ^3x\mathrm{d}x=\int (1-u^2)\mathrm{d}u$.}
\bad{Le changement de variable $u=\sin x$ donne $\displaystyle \int \sin ^3x\mathrm{d}x=\int u^3\mathrm{d}u$.}
\good{Le changement de variable $u=\cos x$ donne $\displaystyle \int \sin ^3x\mathrm{d}x=\int (u^2-1)\mathrm{d}u$.}
\end{answers}
\begin{explanations}
On a : $\cos ^3x=(1-\sin ^2x)\cos x$. Donc, avec $u=\sin x$, on a : $\mathrm{d}u=\cos x\, \mathrm{d}x$ et
$$\int \cos ^3x\mathrm{d}x=\int (1-\sin ^2x)\cos x\mathrm{d}x=\int (1-u^2)\mathrm{d}u.$$
De m\^eme, on a : $\sin ^3x=(1-\cos ^2x)\sin x$. Donc, avec $u=\cos x$, on a : $\mathrm{d}u=-\sin x\mathrm{d}x$ et
$$\int \cos ^3x\mathrm{d}x=\int (1-\cos ^2x)\sin x\mathrm{d}x=-\int (1-u^2)\mathrm{d}u.$$
\end{explanations}
\end{question}
}
